\section*{الفصل \ref{ch:b1:rock-forming-minerals} --- المعادن المكوِّنة للصخور}

\begin{solution}{exo:b1:rock-forming-minerals:1}
\ce{SiO4^{4-}}: \omterm{def:b1:rock-forming-minerals:silicate-classes}{السيليكات الجزيرية}، $s = 0$. \ce{Si2O7^{6-}}: المزدوجة،
$s = 1$. \ce{Si6O18^{12-}}: الحلقية، $s = 2$. \ce{Si2O5^{2-}}:
الصفائحية، $s = 3$. \ce{SiO2}: الإطارية، $s = 4$. \ce{SiO3^{2-}}
له $4 - s/2 = 3$، أي $s = 2$: حلقة (الحلقية) أو سلسلة مفردة
(السلسلية)؛ والنسبة وحدها لا تميّز بينهما.
\end{solution}

\begin{solution}{exo:b1:rock-forming-minerals:2}
الفورستيريت: $2\times24.305 + 28.085 + 4\times15.999 =
\qty{140.691}{g.mol^{-1}}$؛ الفايالايت: $2\times55.845 + 28.085 +
4\times15.999 = \qty{203.771}{g.mol^{-1}}$. الحجوم المولية:
$140.691/3.275 = \qty{42.96}{cm^3.mol^{-1}}$ و$203.771/4.392 =
\qty{46.40}{cm^3.mol^{-1}}$.
\end{solution}

\begin{solution}{exo:b1:rock-forming-minerals:3}
$(0.78 - 0.72)/0.72 = +\qty{8.3}{\percent}$ للحديد \ce{Fe^{2+}}؛
$-\qty{4.2}{\percent}$ للنيكل \ce{Ni^{2+}}؛ $+\qty{15.3}{\percent}$ للمنغنيز \ce{Mn^{2+}}؛ $+\qty{38.9}{\percent}$ للكالسيوم \ce{Ca^{2+}}. الحديد والنيكل،
تحت \qty{15}{\percent}، يحلّان محل المغنيسيوم بحرية؛ والمنغنيز، فوق ذلك بقليل، محدود؛ والكالسيوم، فوق \qty{30}{\percent}، بالكاد يدخل. الأربعة كلها لها الشحنة نفسها التي للمغنيسيوم.
\end{solution}

\begin{solution}{exo:b1:rock-forming-minerals:4}
$(3.51 - 2.16)/2.16 = \qty{62.5}{\percent}$ من الغرافيت إلى الماس؛
$(2.93 - 2.71)/2.71 = \qty{8.1}{\percent}$ من الكالسيت إلى الأراغونيت. والمتعدّدات الأشكال الأكثف، الماس والأراغونيت، هي الأشكال ذات الضغط المرتفع.
\end{solution}

\begin{solution}{exo:b1:rock-forming-minerals:5}
مع $x = 0.30$ في $\mathrm{Na}_{1-x}\mathrm{Ca}_x\mathrm{Al}_{1+x}\mathrm{Si}_{3-x}\mathrm{O}_8$:
$\mathrm{(Na_{0.7}Ca_{0.3})(Al_{1.3}Si_{2.7})O_8}$. الشحنات:
$0.7 + 2\times0.3 + 3\times1.3 + 4\times2.7 = 0.7 + 0.6 + 3.9 + 10.8 = +16$،
و$8\times(-2) = -16$.
\end{solution}

\begin{solution}{exo:b1:rock-forming-minerals:6}
تتكرّر وحدة فيها رباعيان يتشاركان ثلاثة رؤوس (وهما اللذان يصل رأسهما السلسلتين) ورباعيان يتشاركان رأسين. للرباعيات الأربعة $4\times4 = 16$ رأسًا، منها $3 + 3 + 2 + 2 = 10$ مشتركة، أي $10/2 = 5$ أكسجينات مشتركة؛ أما $16 - 10 = 6$ الباقية فغير مشتركة. المجموع $5 + 6 = 11$ أكسجينًا لأربع ذرات سيليكون: \ce{Si4O11}، بشحنة $4\times4 - 11\times2 = -6$. والقضية تعطي النتيجة نفسها ب$s = 5/2$: $4 - 5/4 = 11/4$ أكسجينًا لكل سيليكون، وشحنة $5/2 - 4 = -3/2$ لكل سيليكون.
\end{solution}

\begin{solution}{exo:b1:rock-forming-minerals:7}
عدد المولات في \qty{100}{g}: \ce{SiO2} $40.20/60.083 = 0.66907$؛ \ce{MgO}
$47.10/40.304 = 1.16862$؛ \ce{FeO} $12.20/71.844 = 0.16981$؛ \ce{NiO}
$0.40/74.692 = 0.00536$. الأكسجينات: $2\times0.66907 + 1.16862 + 0.16981 +
0.00536 = 2.68194$؛ $f = 4/2.68194 = 1.49146$. لكل أربعة أكسجينات: Si
\num{0.998}، وMg \num{1.743}، وFe \num{0.253}، وNi \num{0.008}. تحوي المواقع
الثمانية $1.743 + 0.253 + 0.008 = 2.004$، والموقع الرباعي
\num{0.998}: تحليل جيد، \ce{(Mg_{1.743}Fe_{0.253}Ni_{0.008})SiO4}.
$\mathrm{Fo} = 100\times1.743/(1.743 + 0.253) = 87.3$. يضرب العامل $f$ المغنيسيوم والحديد بالقدر نفسه، فيحذف في النسبة:
$\mathrm{Fo} = 100\,(47.10/40.304)/(47.10/40.304 + 12.20/71.844) =
100\times1.16862/1.33843 = 87.3$.
\end{solution}

\begin{solution}{exo:b1:rock-forming-minerals:8}
الموجب: $1 + 2\times3 + 3\times4 + 3 = +22$؛ والسالب:
$10\times(-2) + 2\times(-1) = -22$. طبقة تركيبها \ce{Si4O10} تحمل $-4$ (\cref{prop:b1:rock-forming-minerals:si-o}، $s = 3$)؛ وإحلال سيليكون واحد من كل أربعة بالألومنيوم يضيف شحنة سالبة واحدة، ف\ce{(Si3Al)O10} تحمل $-5$. الألومنيومان الثمانيان مع الهيدروكسيلين، $6 - 2 = +4$، يوازنان أربعة من الخمسة؛ وأيون البوتاسيوم بين الصفائح يوازن الخامسة.
\end{solution}

\begin{solution}{exo:b1:rock-forming-minerals:9}
(أ) $M = (140.691 + 203.771)/2 = \qty{172.231}{g.mol^{-1}}$، $V = (42.96 +
46.40)/2 = \qty{44.68}{cm^3.mol^{-1}}$، $\rho = 172.231/44.68 =
\qty{3.855}{g.cm^{-3}}$؛ ومتوسط الكثافتين $(3.275 + 4.392)/2 =
\qty{3.834}{g.cm^{-3}}$، أي \qty{0.021}{g.cm^{-3}} أقل. (ب) نحلّ
$203.771 + x(140.691 - 203.771) = 3.600\,\bigl(46.40 + x(42.96 - 46.40)\bigr)$:
$203.771 - 63.080x = 167.040 - 12.384x$، ف$x = 36.731/50.696 = 0.7245$: \omterm{def:b1:rock-forming-minerals:groups}{الأوليفين} هو $\text{Fo}_{72}$ (إلى أقرب نسبة مئوية).
\end{solution}

\begin{solution}{exo:b1:rock-forming-minerals:10}
بما أن $M_i = \rho_i V_i$،
\[
  \rho(x) - \bigl(x\rho_1 + (1 - x)\rho_2\bigr)
  = \frac{x\rho_1V_1 + (1 - x)\rho_2V_2 - \bigl(x\rho_1 + (1 - x)\rho_2\bigr)\bigl(xV_1 + (1 - x)V_2\bigr)}{V(x)} .
\]
بفك الجداء، يكون البسط $x\rho_1V_1 - x^2\rho_1V_1 +
(1 - x)\rho_2V_2 - (1 - x)^2\rho_2V_2 - x(1 - x)(\rho_1V_2 + \rho_2V_1)$، و
$x - x^2 = x(1 - x)$، و$(1 - x) - (1 - x)^2 = x(1 - x)$، فيكون
$x(1 - x)(\rho_1V_1 + \rho_2V_2 - \rho_1V_2 - \rho_2V_1) =
x(1 - x)(\rho_1 - \rho_2)(V_1 - V_2)$. بالنسبة إلى \omterm{def:b1:rock-forming-minerals:groups}{الأوليفين} (1 = الفورستيريت) الفرقان سالبان، $3.275 - 4.392 = -1.117$ و$42.96 - 46.40 =
-3.44$، فجداؤهما موجب، وتقع الكثافة فوق الخط المستقيم. عند $\text{Fo}_{50}$: $0.25\times1.117\times3.44/44.68 =
\qty{0.0215}{g.cm^{-3}}$، وهذا يتفق مع \cref{exo:b1:rock-forming-minerals:9} في حدود تقريب الحجوم المولية.
\end{solution}

\begin{solution}{exo:b1:rock-forming-minerals:11}
المجموع \qty{99.60}{\percent}. مولات الأكاسيد: \ce{SiO2} $0.90042$،
\ce{Al2O3} $0.28344$، \ce{CaO} $0.19973$، \ce{Na2O} $0.08229$،
\ce{K2O} $0.00318$. الأكسجينات: $1.80084 + 0.85033 + 0.19973 + 0.08229 +
0.00318 = 2.93636$؛ $f = 8/2.93636 = 2.72446$. الكاتيونات (ضعف المولات في
\ce{Al2O3} و\ce{Na2O} و\ce{K2O}): Si \num{2.453}، وAl \num{1.544}، وCa
\num{0.544}، وNa \num{0.448}، وK \num{0.017}. الرباعيات: $2.453 + 1.544 = 3.997$؛ والموقع الكبير: $0.544 + 0.448 + 0.017 = 1.009$. ثم
$\mathrm{Al} - 1 = 0.544$، وهو يساوي الكالسيوم في حدود التقريب: كل كالسيوم دخل مع ألومنيوم واحد. $\mathrm{An} = 100\times0.544/1.009 = 53.9$: فهو اللابرادوريت، $\text{An}_{54}$.
\end{solution}

\begin{solution}{exo:b1:rock-forming-minerals:12}
(أ) $M(\ce{Fe3O4}) = 3\times55.845 + 4\times15.999 = \qty{231.531}{g.mol^{-1}}$.
حديده الثلاثة، مُبلَّغًا عنها على صورة \ce{FeO}، تزن $3\times71.844 = \qty{215.532}{g}$،
أي $215.532/231.531 = \qty{93.09}{\percent}$؛ أما الأكسجين المرتبط ب\ce{Fe^{3+}}، الزائد على ما يحسبه \ce{FeO}، فهو غائب عن المجموع. (ب) $93.09/71.844 = 1.29572$ مول من \ce{FeO} ومن الأكسجين؛ $f =
4/1.29572 = 3.08708$؛ أربع ذرات حديد لكل أربعة أكسجينات، \ce{Fe4O4}. والمغنتيت فيه ثلاثة كاتيونات لأربعة أكسجينات. وقد عُدّت الأكسجينات أقل من حقيقتها: كل \ce{Fe^{3+}} يحمل أكسجينًا ونصفًا، لا أكسجينًا واحدًا كما يفترض \ce{FeO}. (ج) على ثلاثة كاتيونات، Fe $= 3$
مع $a$ من \ce{Fe^{2+}} و$b$ من \ce{Fe^{3+}}: $a + b = 3$ و$2a + 3b = 8$، ف$b = 2$، و$a = 1$: $\mathrm{Fe^{2+}Fe^{3+}_2O_4}$. (د) واحد من \ce{FeO} وواحد من \ce{Fe2O3} لكل صيغة وحدوية: $71.844/231.531 = \qty{31.03}{\percent}$
من \ce{FeO} و$159.687/231.531 = \qty{68.97}{\percent}$ من \ce{Fe2O3}، والمجموع \qty{100.00}{\percent}.
\end{solution}

\begin{solution}{pb:b1:rock-forming-minerals:1}
\textbf{1.} \ce{SiO2} $28.085 + 2\times15.999 = 60.083$؛ \ce{FeO} 71.844؛
\ce{MnO} 70.937؛ \ce{MgO} 40.304؛ \ce{NiO} 74.692؛ \ce{CaO} 56.077
(\si{g.mol^{-1}}).

\textbf{2.} $39.50 + 16.53 + 0.23 + 43.32 + 0.20 + 0.22 =
\qty{100.00}{\percent}$. ولو بلغ المجموع \qty{97}{\percent} لأوحى ذلك بمكوّن لم يُحلَّل (ماء، أو أكسجين مرتبط ب\ce{Fe^{3+}} مُبلَّغ عنه على صورة \ce{FeO})، أو بتحليل رديء.

\textbf{3.} \ce{SiO2} $39.50/60.083 = 0.65742$؛ \ce{FeO} 0.23008؛ \ce{MnO}
0.00324؛ \ce{MgO} 1.07483؛ \ce{NiO} 0.00268؛ \ce{CaO} 0.00392.

\textbf{4.} $2\times0.65742 = 1.31485$ ل\ce{SiO2}، والأعداد نفسها من السؤال 3 للباقي؛ والمجموع $1.31485 + 0.23008 + 0.00324 + 1.07483 +
0.00268 + 0.00392 = 2.62960$.

\textbf{5.} لكل أكسيد كاتيون واحد: الأعداد نفسها من السؤال 3.

\textbf{6.} $f = 4/2.62960 = 1.52114$.

\textbf{7.} Si $0.65742\times1.52114 = 1.000$؛ Fe 0.350؛ Mn 0.005؛ Mg 1.635؛
Ni 0.004؛ Ca 0.006.

\textbf{8.} $1.635 + 0.350 + 0.005 + 0.004 + 0.006 = 2.000$ للمواقع الثمانية، و\num{1.000} للموقع الرباعي: وهي تماما القيم المثالية 2 و\num{1}.

\textbf{9.} \ce{(Mg_{1.635}Fe_{0.350}Mn_{0.005}Ni_{0.004}Ca_{0.006})SiO4}.

\textbf{10.} $2\times2.000 + 4\times1.000 = +8.000$ مقابل $4\times(-2) =
-8$. كل كاتيون حُسب بشحنة أكسيده، فالشحنات الموجبة تساوي ضعف الأكسجينات بالبناء: لذلك لا يمكن للتحقق أن يفشل. أما مجاميع المواقع في السؤال 8 فكان يمكن أن تفشل.

\textbf{11.} يختلف النيكل عن المغنيسيوم بنسبة $(0.69 - 0.72)/0.72 = -\qty{4}{\percent}$، وهي دون \qty{15}{\percent} بكثير: فهو يحل محله بحرية. أما الكالسيوم فأكبر منه بنسبة $(1.00 - 0.72)/0.72 = \qty{39}{\percent}$، أي فوق \qty{30}{\percent}: فهو بالكاد يدخل.

\textbf{12.} الحديد المحسوب على صورة \ce{FeO} حين يكون جزء منه \ce{Fe^{3+}} يعطي أكسجينات أقل من اللازم لكل كاتيون، ومن ثمّ أكثر من ثلاثة كاتيونات لأربعة أكسجينات (\cref{exo:b1:rock-forming-minerals:12}). المجموع $1.000 + 2.000 = 3.000$: لا يوجد \ce{Fe^{3+}} يُعتد به.

\textbf{13.} $\mathrm{Fo} = 100\times1.635/(1.635 + 0.350) = 163.5/1.985 =
82.4$.

\textbf{14.} $\mathrm{Mg\#} = 82.4$: كل الحديد \ce{Fe^{2+}} (السؤال 12)، فإن $\mathrm{Fe^{2+}} = \mathrm{Fe}$، فالتعريفان يتطابقان.

\textbf{15.} $\text{Fo}_{82.4}\text{Fa}_{17.6}$، $x = 0.824$؛ يملأ الحديد \qty{17.6}{\percent} من المواقع الثمانية في النموذج الثنائي ($0.350/2.000 = \qty{17.5}{\percent}$ مع العناصر الثانوية).

\textbf{16.} يضرب العامل المغنيسيوم والحديد بالقدر نفسه فيحذف: $\mathrm{Fo} = 100\times1.07483/(1.07483 + 0.23008) = 82.4$.

\textbf{17.} $40.5/40.304 = 1.00486$ و$20.0/71.844 = 0.27838$؛
$\mathrm{Fo} = 100\times1.00486/1.28324 = 78.3$. \omterm{def:b1:rock-forming-minerals:groups}{فالأوليفين} المتأخر هو الأغنى بالحديد.

\textbf{18.} $M = 140.691$ و\qty{203.771}{g.mol^{-1}}؛
$V = 140.691/3.275 = 42.96$ و$203.771/4.392 =
\qty{46.40}{cm^3.mol^{-1}}$.

\textbf{19.} $M = 0.824\times140.691 + 0.176\times203.771 = 115.929 +
35.864 = \qty{151.793}{g.mol^{-1}}$؛ $V = 0.824\times42.96 +
0.176\times46.40 = 35.399 + 8.166 = \qty{43.565}{cm^3.mol^{-1}}$.

\textbf{20.} $\rho = 151.793/43.565 = \qty{3.484}{g.cm^{-3}}$.

\textbf{21.} $0.824\times3.275 + 0.176\times4.392 = 2.699 + 0.773 =
\qty{3.472}{g.cm^{-3}}$، أي \qty{0.012}{g.cm^{-3}} (\qty{0.3}{\percent}) أدنى من قيمة المزج المثالي: قريبة، كما يتنبأ \cref{exo:b1:rock-forming-minerals:10}.

\textbf{22.} $0.12\times\qty{3484}{kg.m^{-3}}\times\qty{1}{m^3} =
\qty{418}{kg}$.

\textbf{23.} $418/2950 = \qty{14.2}{\percent}$: يشكّل \omterm{def:b1:rock-forming-minerals:groups}{الأوليفين} الكثيف حصة من الكتلة أكبر من حصته في الحجم.

\textbf{24.} \omterm{def:b1:rock-forming-minerals:groups}{الأوليفين} هو $\text{Fo}_{82}$: محتواه من الفورستيريت \textbf{82 في المئة}.
\end{solution}
